% GENERATED FILE. Edit canonical lessons, then run npm run generate:books.
% canonical-source-hash: fnv1a64:ccf059d64f544666
% canonical-lessons: TA-C363-pullivivaram, TA-C363-catavitam, TA-C363-caracari, TA-C363-vikitam, TA-C363-kani

\chapter{Numbers in an Argument}
\label{ch:ta-numbers-in-an-argument}
\hlchaptermodality{363}

\begin{chapteropening}
\textbf{By the end of this chapter:} \emph{I can use statistics, percentages, averages and ratios.}

Complete the last lesson of chapter 363: I can use statistics, percentages, averages and ratios.
\end{chapteropening}

\section[\texorpdfstring{\ta{புள்ளிவிவரம்} (puḷḷivivaram)}{புள்ளிவிவரம் (puḷḷivivaram)}]{\ta{புள்ளிவிவரம்} (puḷḷivivaram) — statistics}

\label{lesson:TA-C363-pullivivaram}



\begin{quote}
Before the new one: say the Tamil for an analysis, then the Tamil for to analyse closely.
\end{quote}

\subsection*{You'll want to know: \ta{புள்ளிவிவரம்}}
\textbf{\ta{புள்ளிவிவரம்}} — \emph{puḷḷivivaram} — \textquotedblleft{}statistics\textquotedblright{}.

\begin{quote}
\ta{இந்தப்} \ta{புள்ளிவிவரங்கள்} \ta{எதைக்} \ta{காட்டுகின்றன}?
\end{quote}

\begin{quote}
\emph{intap puḷḷivivaraṅkaḷ etaik kāṭṭukiṉṟaṉa?}
\end{quote}

\begin{quote}
What do these statistics show?
\end{quote}

Usually plural: \ta{புள்ளிவிவரங்கள்}.

\begin{cousinweb}
\textbf{\ta{புள்ளிவிவரம்}} (\emph{puḷḷivivaram}) joins native \emph{puḷḷi}, \textquotedblleft{}mark, dot\textquotedblright{} (DEDR 4327), and \emph{vivaram}, \textquotedblleft{}details\textquotedblright{}, borrowed from Sanskrit \emph{vivaraṇa}, \textquotedblleft{}laying open, explanation\textquotedblright{} (Monier-Williams): details by the point.
\end{cousinweb}

\begin{grammarlens}[title={What you've built}]
One more word for giving reasons and opinions.
\end{grammarlens}

\subsection*{Your turn}
\begin{itemize}
\raggedright
  \item \emph{Say it:} \emph{puḷḷivivaram}
  \item \emph{Say it:} \emph{puḷḷivivaram}, once more
  \item \emph{Say it:} \emph{intap puḷḷivivaraṅkaḷ etaik kāṭṭukiṉṟaṉa?} — What do these statistics show?
  \item \emph{Recall:} say the Tamil for an analysis, then the Tamil for to analyse closely, then say \emph{puḷḷivivaram} again
\end{itemize}

\begin{tcolorbox}[breakable,colback=teal!4,colframe=teal!35!black,title={Before you move on}]
What does \ta{புள்ளிவிவரம்} mean? (\textquotedblleft{}Statistics\textquotedblright{}.) Say it once more.
\end{tcolorbox}
\section[\texorpdfstring{\ta{சதவீதம்} (catavītam)}{சதவீதம் (catavītam)}]{\ta{சதவீதம்} (catavītam) — a percentage; per cent}

\label{lesson:TA-C363-catavitam}



\begin{quote}
Before the new one: say the Tamil for to analyse closely, then the Tamil for statistics.
\end{quote}

\subsection*{You'll want to know: \ta{சதவீதம்}}
\textbf{\ta{சதவீதம்}} — \emph{catavītam} — \textquotedblleft{}a percentage; per cent\textquotedblright{}.

\begin{quote}
\ta{எண்பது} \ta{சதவீதம்} \ta{மாணவர்கள்} \ta{தேர்வில்} \ta{தேர்ச்சி} \ta{பெற்றார்கள்}.
\end{quote}

\begin{quote}
\emph{eṇpatu catavītam māṇavarkaḷ tērvil tērcci peṟṟārkaḷ.}
\end{quote}

\begin{quote}
Eighty per cent of the students passed the exam.
\end{quote}

Numbers are spelled out: \ta{பத்து} \ta{சதவீதம்}, ten per cent.

\begin{cousinweb}
\textbf{\ta{சதவீதம்}} (\emph{catavītam}) joins \emph{catam}, \textquotedblleft{}a hundred\textquotedblright{}, borrowed from Sanskrit \emph{śata} (CDIAL 12278), and \emph{vītam}, \textquotedblleft{}a rate, a share\textquotedblright{}, whose origin is not settled: so much per hundred.
\end{cousinweb}

\begin{grammarlens}[title={What you've built}]
One more word for giving reasons and opinions.
\end{grammarlens}

\subsection*{Your turn}
\begin{itemize}
\raggedright
  \item \emph{Say it:} \emph{catavītam}
  \item \emph{Say it:} \emph{catavītam}, once more
  \item \emph{Say it:} \emph{eṇpatu catavītam māṇavarkaḷ tērvil tērcci peṟṟārkaḷ.} — Eighty per cent of the students passed the exam.
  \item \emph{Recall:} say the Tamil for to analyse closely, then the Tamil for statistics, then say \emph{catavītam} again
\end{itemize}

\begin{tcolorbox}[breakable,colback=teal!4,colframe=teal!35!black,title={Before you move on}]
What does \ta{சதவீதம்} mean? (\textquotedblleft{}A percentage; per cent\textquotedblright{}.) Say it once more.
\end{tcolorbox}
\section[\texorpdfstring{\ta{சராசரி} (carācari)}{சராசரி (carācari)}]{\ta{சராசரி} (carācari) — an average}

\label{lesson:TA-C363-caracari}



\begin{quote}
Before the new one: say the Tamil for statistics, then the Tamil for a percentage.
\end{quote}

\subsection*{You'll want to know: \ta{சராசரி}}
\textbf{\ta{சராசரி}} — \emph{carācari} — \textquotedblleft{}an average\textquotedblright{}.

\begin{quote}
\ta{இந்த} \ta{வகுப்பின்} \ta{சராசரி} \ta{மதிப்பெண்} \ta{எழுபது}.
\end{quote}

\begin{quote}
\emph{inta vakuppiṉ carācari matippeṇ eḻupatu.}
\end{quote}

\begin{quote}
The average mark in this class is seventy.
\end{quote}

\ta{சராசரியாக}, on average.

\begin{cousinweb}
\textbf{\ta{சராசரி}} (\emph{carācari}) is usually traced to Persian \emph{sarāsar}, \textquotedblleft{}from end to end, throughout\textquotedblright{}, which reached Tamil through Hindustani: a figure taken across the whole.
\end{cousinweb}

\begin{grammarlens}[title={What you've built}]
One more word for giving reasons and opinions.
\end{grammarlens}

\subsection*{Your turn}
\begin{itemize}
\raggedright
  \item \emph{Say it:} \emph{carācari}
  \item \emph{Say it:} \emph{carācari}, once more
  \item \emph{Say it:} \emph{inta vakuppiṉ carācari matippeṇ eḻupatu.} — The average mark in this class is seventy.
  \item \emph{Recall:} say the Tamil for statistics, then the Tamil for a percentage, then say \emph{carācari} again
\end{itemize}

\begin{tcolorbox}[breakable,colback=teal!4,colframe=teal!35!black,title={Before you move on}]
What does \ta{சராசரி} mean? (\textquotedblleft{}An average\textquotedblright{}.) Say it once more.
\end{tcolorbox}
\section[\texorpdfstring{\ta{விகிதம்} (vikitam)}{விகிதம் (vikitam)}]{\ta{விகிதம்} (vikitam) — a ratio; a rate}

\label{lesson:TA-C363-vikitam}



\begin{quote}
Before the new one: say the Tamil for a percentage, then the Tamil for an average.
\end{quote}

\subsection*{You'll want to know: \ta{விகிதம்}}
\textbf{\ta{விகிதம்}} — \emph{vikitam} — \textquotedblleft{}a ratio; a rate\textquotedblright{}.

\begin{quote}
\ta{இந்த} \ta{ஊரில்} \ta{ஆண்} \ta{பெண்} \ta{விகிதம்} \ta{சமமாக} \ta{இருக்கிறது}.
\end{quote}

\begin{quote}
\emph{inta ūril āṇ peṇ vikitam camamāka irukkiṟatu.}
\end{quote}

\begin{quote}
In this town the ratio of men to women is equal.
\end{quote}

\ta{என்ற} \ta{விகிதத்தில்}, in the ratio of: \ta{ஒன்றுக்கு} \ta{இரண்டு} \ta{என்ற} \ta{விகிதத்தில்}, one to two.

\begin{cousinweb}
\textbf{\ta{விகிதம்}} (\emph{vikitam}) is borrowed from Sanskrit \emph{vihita}, \textquotedblleft{}distributed, apportioned, fixed\textquotedblright{} (Monier-Williams). Tamil writes the Sanskrit h between vowels as \ta{க}.
\end{cousinweb}

\begin{grammarlens}[title={What you've built}]
One more word for giving reasons and opinions.
\end{grammarlens}

\subsection*{Your turn}
\begin{itemize}
\raggedright
  \item \emph{Say it:} \emph{vikitam}
  \item \emph{Say it:} \emph{vikitam}, once more
  \item \emph{Say it:} \emph{inta ūril āṇ peṇ vikitam camamāka irukkiṟatu.} — In this town the ratio of men to women is equal.
  \item \emph{Recall:} say the Tamil for a percentage, then the Tamil for an average, then say \emph{vikitam} again
\end{itemize}

\begin{tcolorbox}[breakable,colback=teal!4,colframe=teal!35!black,title={Before you move on}]
What does \ta{விகிதம்} mean? (\textquotedblleft{}A ratio; a rate\textquotedblright{}.) Say it once more.
\end{tcolorbox}
\section[\texorpdfstring{\ta{கணி} (kaṇi)}{கணி (kaṇi)}]{\ta{கணி} (kaṇi) — to predict; to calculate}

\label{lesson:TA-C363-kani}



\begin{quote}
Before the new one: say the Tamil for an average, then the Tamil for a ratio.
\end{quote}

\subsection*{You'll want to know: \ta{கணி}}
\textbf{\ta{கணி}} — \emph{kaṇi} — \textquotedblleft{}to predict; to calculate\textquotedblright{}.

\begin{quote}
\ta{அடுத்த} \ta{ஆண்டு} \ta{விலைகள்} \ta{உயரும்} \ta{என்று} \ta{அவர்கள்} \ta{கணித்தார்கள்}.
\end{quote}

\begin{quote}
\emph{aṭutta āṇṭu vilaikaḷ uyarum eṉṟu avarkaḷ kaṇittārkaḷ.}
\end{quote}

\begin{quote}
They predicted that prices would rise next year.
\end{quote}

Past (I …): \ta{கணித்தேன்}, \emph{kaṇittēṉ}.

The prediction is a clause with \ta{என்று}. \ta{கணிதம்}, mathematics, comes from the same Sanskrit root.

\begin{cousinweb}
\textbf{\ta{கணி}} (\emph{kaṇi}) is a verb built on Sanskrit \emph{gaṇayati}, \textquotedblleft{}counts, takes notice of\textquotedblright{} (CDIAL 3993), the source of \ta{கணிதம்}, \textquotedblleft{}mathematics\textquotedblright{}. Tamil writes the Sanskrit g as \ta{க}.
\end{cousinweb}

\begin{grammarlens}[title={What you've built}]
One more word for giving reasons and opinions.
\end{grammarlens}

\subsection*{Your turn}
\begin{itemize}
\raggedright
  \item \emph{Say it:} \emph{kaṇi}
  \item \emph{Say it:} \emph{kaṇi}, once more
  \item \emph{Say it:} \emph{aṭutta āṇṭu vilaikaḷ uyarum eṉṟu avarkaḷ kaṇittārkaḷ.} — They predicted that prices would rise next year.
  \item \emph{Recall:} say the Tamil for an average, then the Tamil for a ratio, then say \emph{kaṇi} again
\end{itemize}

\begin{tcolorbox}[breakable,colback=teal!4,colframe=teal!35!black,title={Before you move on}]
What does \ta{கணி} mean? (\textquotedblleft{}To predict; to calculate\textquotedblright{}.) Say it once more.
\end{tcolorbox}
